% Source: https://artofproblemsolving.com/community/c3h1993312_i_might_have_asked_this_before
% Stable identity: aops-topic-1993312-q01
A chord of length $r$ together with the two radii to its ends forms an equilateral triangle, so it cuts off an arc of $60^\circ$. ``Drawn in a clockwise direction'' means: from the chosen point $X$, the chord goes to the point $X'$ that is $60^\circ$ clockwise from $X$. So each chord is determined by its starting point and spans the $60^\circ$ arc just clockwise of it.

Let the first point be $X$ and measure the position of the second point $Y$ by the counterclockwise angle $\theta$ from $X$, uniform on $[0^\circ,360^\circ)$. Two chords cross exactly when their endpoints alternate around the circle, which here means the $60^\circ$ arcs $X'X$ and $Y'Y$ overlap without coinciding. This happens exactly when $Y$ is within $60^\circ$ of $X$ on either side ($0<\theta<60^\circ$ or $300^\circ<\theta<360^\circ$); otherwise the arcs, and hence the chords, are disjoint.

\begin{asy}
usepackage("amsmath");
size(4.5cm);
defaultpen(fontsize(10pt));
pair X=dir(90), X2=dir(30), Y=dir(120), Y2=dir(60);
draw(unitcircle,linewidth(0.9));
draw(X--X2,blue+linewidth(1.1)); draw(Y--Y2,red+linewidth(1.1));
draw(arc((0,0),1.32,84,36),blue,EndArrow(TeXHead));
dot(X,blue);dot(X2,blue);dot(Y,red);dot(Y2,red);dot((0,0));
label("$X$",X,(-0.5,1),blue); label("$X'$",X2,E,blue); label("$Y$",Y,NW,red); label("$Y'$",Y2,NE,red);
draw((0,0)--X,dashed+gray(0.5)); draw((0,0)--X2,dashed+gray(0.5));
label("$60^\circ$",(0.22,0.6),gray(0.3));
label("\mbox{}",(min(currentpicture,user=true)+max(currentpicture,user=true))/2,black); // keeps the text after the figure black when LaTeX typesets the labels inline
\end{asy}

Hence the probability is $\dfrac{120}{360}=\boxed{\dfrac{1}{3}}$, choice $\textbf{(D)}$.
